\documentclass[12pt,addpoints,answers]{exam}
\usepackage[utf8]{inputenc}
\usepackage{amsmath,amsfonts,amssymb,amsthm}
\usepackage[margin=1in]{geometry}
\usepackage{mathtools}
\usepackage{hyperref}
\usepackage{fullpage}
\usepackage{microtype}
\usepackage{xspace}
\usepackage[svgnames]{xcolor}
\usepackage[sc]{mathpazo}
\usepackage{enumitem}
\usepackage{bm}
\usepackage{algorithm}
\usepackage{algpseudocode}

\pagestyle{head}

%----------Header--------------------%
\def\course{{\sc Foundations and Applications of Blockchains}}
\def\term{Depto. de Computaci\'{o}n, $2^{\textrm{do}}$ Cuatrimestre 2023}
\def\prof{Lecturer: Juan Garay}
\newcommand{\handout}[5]{
	\renewcommand{\thepage}{\arabic{page}}
	\begin{center}
		\framebox{
			\vbox{
				\hbox to 5.78in { \hfill \large{\course} \hfill }
				\vspace{2mm}
				%    \hbox to 5.78in { \hfill \large{\prof} \hfill }
				%       \vspace{2mm}
				\hbox to 5.78in { {\Large \hfill \textbf{#5}  \hfill} }
				\vspace{2mm}
				\hbox to 5.78in { \term \hfill \emph{#2}}
				\hbox to 5.78in { {#3 \hfill \emph{#4}}}
			}
		}
	\end{center}
	\vspace*{4mm}
}
\newcommand{\hw}[4]{\handout{#1}{#2}{#3}{#4}{Homework #1}}

% -- For ignoring stuff -- %
\newcommand{\ignore}[1]{}

\begin{document}
	
	%----Specs: change accordingly-----%
	\newif\ifstudent % comment out false
	\studenttrue 
	% \studentfalse
	
	\def\hwnum{1}
	\def\issuedate{24/8/23}
	\def\duedate{25/9/23, 3:30pm} % 
	\def\yourname{} % put your name here
%------------------------------%

\ifstudent
\hw{\hwnum}{\issuedate}{Student: \yourname}{Due: \duedate}%
\else
\hw{\hwnum}{\issuedate}{\prof}{Due: \duedate}%
\fi

% \ignore{

\noindent \textbf{Instructions}

\begin{itemize}
    \item The way to deliver your solution will be explained in the following classes.   
      
     If you are proficient with \LaTeX, you may also typeset your submission and submit in PDF format. To do so, uncomment the ``\%\textbackslash begin\{solution\}'' and ``\%\textbackslash end\{solution\}'' lines and write your solution between those two command lines.
    
      \item Your solutions will be graded on \emph{correctness} and
    \emph{clarity}. You should only submit work that you believe to be
    correct.
    % , and you will get significantly more partial credit if you     clearly identify the gap(s) in your solution.
    
    \item You may collaborate with others on Questions 1-3 of this problem set.  However,
    you must \textbf{{write up your own solutions}} and \textbf{{list
      your collaborators and any external sources}} for each
    problem. Be ready to explain your solutions orally to a member of the course staff
    if asked. You {\bf may not} collaborate with others on Question 4.
    
    \ignore{
    \item For problems that require you to provide an algorithm, you must
    give a precise description of the algorithm, together with a proof
    of correctness and an analysis of its running time. You may use
    algorithms from class as subroutines. You may also use any facts
    that we proved in class or from the book.
    } %IGNORE
    
\end{itemize}

\noindent This homework contains \numquestions\ questions,
% \numpages\ pages
for a total of \numpoints \ points.
% and \numbonuspoints\ bonus points.

%\medskip
\newpage

%} %IGNORE

\begin{questions}


\question We saw in class the notion of {\em digital signatures} and their security properties, {\em existential unforgeability} being an important one. 

\begin{parts}
    \part[5] Describe the purpose of a {\em Public-Key Infrastructure} (PKI). Can the security properties of a digital signature be guaranteed without a PKI? Elaborate.
    
    \begin{solution} %Uncomment and type your solution here
      El propósito de PKI es proveer a los distintos nodos de una red una manera de verificar la autenticidad de los mensajes enviados por otros participantes.
      En particular, en clase vimos el problema de los generales bizantinos (la version de que uno le envia a todos un mensaje).
      Si ese mensaje estuviese firmado, podriamos verificar que viene realmente de ese "general bizantino". 
      Pero para verificar la clave pública de ese general, necesitamos que esté autorizado por una CA (\textit{Certificate Authority}).

      Las propiedades de una firma digital son las siguientes:
      \begin{itemize}
        \item Autenticidad: el mensaje fue enviado por quien lo firmó.
        \item Integridad: el mensaje no fue modificado.
        \item No repudio: el firmante no puede negar que lo firmó.
        \item \textit{Existential unforgeability}: un atacante no puede crear otra firma válida aunque conozca otras firmas de ese mismo firmante.
      \end{itemize}

      Para cumplir la autenticidad sin PKI, necesitaríamos que cada nodo comparta con los demás su clave pública de una manera confiable.
      Esto es para que, cuando reciba una firma, pueda verificar que coincide con la verdadera clave pública del firmante. Esto asumiría que la clave pública que tengo de ese participante es la correcta.
      Con PKI sería mas simple, ya que cada nodo puede verificar la clave pública de otro nodo a través de la CA.

      El no-repudio es particularmente problemático, ya que el firmante podría simplemente alegar que esa no es su clave pública. Si tuviésemos PKI, esta entidad podría intervenir y decir que esa no es la clave pública del firmante.
      La integridad y \textit{existential unforgeability} se pueden cumplir sin PKI, ya que dependen del algoritmo de firma en sí mismo. 
    \end{solution}

    \part[5] Bitcoin transactions use the ECDSA signature scheme. Does Bitcoin assume a PKI? If not, reconcile with the above argument.
 
    \begin{solution} %Uncomment and type your solution here
      Bitcoin no asume un PKI. La razón es porque cada participante de una transacción se identifica por su "dirección", que es justamente su clave pública.
      Por ende, el hecho de el mensaje esté firmado por esa dirección nos permite verificar que fue enviado por el dueño.

      Si la firma no es válida, la transacción se va a rechazar por los nodos de la red. Esto garantiza que solamente quien tiene la clave privada puede realizar transacciones.
      En relación a lo dicho anteriormente, vemos que el no repudio se cumple en el sentido de que quien envió una transacción no puede negar haber tenido la clave privada de esa dirección.
      En cuanto a la autenticación de que una transacción fue enviada por el dueño de la dirección, vemos que esto se cumple por lo explicado anteriormente: la clave pública es justamente la dirección de quien envía los fondos.
    \end{solution}
        
\end{parts}

\newpage

\question[10] Refer to Algorithm 4 (Main Loop) in [GKL15]. We saw in class that blockchains would start from a ``genesis'' block, which would provide an unpredictable string to start mining from. Yet, notice that in Algorithm 4 mining starts from the empty string (${\cal C} \leftarrow \varepsilon$). How come? Explain why this works.

\begin{solution} %Uncomment and type your solution here
  Vemos que al inicio del algoritmo (${\cal C} \leftarrow \varepsilon$). En la ronda número 2, vale que $init = False$, entonces se entra a la rama else del if.

  Luego, se busca la máxima cadena válida entre las disponibles, incluyendo ${\cal C} = \varepsilon$. Asumimos que no hay otras disponibles en $RECEIVE()$, 
  entonces que ${\tilde {\cal C}} = \varepsilon$, pues es válida.

  Vemos que es válida observando el código de $validate$, en el que se retorna $ b = V(\textbf{x}_{\cal C}$). Por convención (dicho en el paper), vemos que $\textbf{x}_{\cal C}$ de una cadena vacía es siempre $1$.
  Por ende, la cadena vacía es válida.

  Luego, se genera el valor $x$ con la funcion $I(\cdot)$ y ${\cal C}_{new}$ con $pow$.
  Vemos que esta ${\cal C}_{new}$ generada va a ser el resultado de que se tome $s = 0$ en $pow$, y que se agregue un
  nuevo bloque a ${\tilde {\cal C}}$. Este nuevo bloque va a ser $\langle s, x, ctr \rangle$, con $s = 0$.
  El $x$ y $ctr$ van a ser los que cumplan con $H(ctr, h) < T$, siendo $h = G(0, x)$

  Luego, obtuvimos una nueva cadena (${\cal C} \leftarrow {\cal C}_{new}$) con dos bloques, el vacío y el que generamos recién.
  Las rondas subsiguientes entrarían en los casos en los que se van a usar otras cadenas (no la vacía) como ${\tilde{\cal C}}$, ya que
  ${\cal C}$ contendrá un bloque válido.

  
\end{solution}

\newpage

\question Refer to Algorithm 1 ({\sf validate}) in [GKL15], which implements the {\em chain validation predicate}. 

\begin{parts}

\part[5] Rewrite {\sf validate} so that it starts checking from the {\em beginning} of the chain. 

\begin{solution} %Uncomment and type your solution here

Tomé en cuenta que, en cada bloque, $s$ representa el hash del bloque anterior. Por eso uso previousHash, que nota el hash del bloque anterior para ser verificado con el siguiente.
Otra cosa a tener en cuenta es que usé $\text{first}$ para obtener el primer bloque de la cadena, y ${\cal C}^{1 \rceil }$ para obtener toda la cadena excepto el primer bloque.


\begin{algorithmic}[1]
\Function{validate}{$\mathcal{C}$}
  \State b $\leftarrow$ V($x_\mathcal{C}$) 
  \If{$b \land \mathcal{C} \neq \varepsilon$}
    \State $\langle s, x, ctr \rangle \leftarrow \text{first}({\cal C})$ 
    \If{NOT ($validblock_{q}^{T}(\langle s, x, ctr \rangle)$)} 
        \State b $\leftarrow$ False
    \Else
        \State previousHash $\leftarrow H(ctr, G(s, x))$ 
        \State $\mathcal{C} \leftarrow {\cal C}^{1 \rceil }$
        
        \While{$\mathcal{C} \neq \varepsilon \land b = \text{True}$}
            \State $\langle s, x, ctr \rangle \leftarrow \text{first}({\cal C})$ \
            
            \If{$validblock_{q}^{T}(\langle s, x, ctr \rangle) \land (s = \text{previousHash})$}
                \State previousHash $\leftarrow H(ctr, G(s, x))$ 
                \State $\mathcal{C} \leftarrow {\cal C}^{1 \rceil }$ 
            \Else
                \State b $\leftarrow$ False 
            \EndIf
        \EndWhile
    \EndIf
  \EndIf 
  \State \Return b
\EndFunction
\end{algorithmic}


\end{solution}

\part[5] Discuss pros and cons of this approach, compared to Algorithm 1's.

\begin{solution} %Uncomment and type your solution here
  La principal desventaja de este algoritmo es que, si queremos chequear la validez de una cadena muy larga, y que probablemente si hay un error, este en los últimos bloques, vamos a tardar mucho hasta llegar hasta ahí, ya que estamos mirando desde adelante hacia atrás.
  Decimos que es más probable que haya un error en los últimos bloques ya que, en el contexto de una blockchain real, los bloques se agregan al final, y los nodos comparten una gran parte de la cadena inicial (que si los nodos son honestos debería ser válida). Por lo tanto, tiene sentido comenzar a validar desde el final.
 

  Otra desventaja es que es un poco más difícil de entender. Como las cadenas están formadas por bloques que apuntan a su bloque anterior, sería intuitivo ir desde adelante hacia atrás.
  En este caso estamos yendo desde atrás hacia adelante. 

  Una ventaja que tiene el algoritmo que va hacia adelante es que si, por alguna razón hubiera errores en los primeros bloques, los detectaríamos antes.
  También puede ser mas fácil de entender el hecho de avanzar hacia adelante, pero una vez que se comprende la estructura de las cadenas y los bloques, queda claro que ir hacia atrás es hacer mejor uso de las estructuras de datos.


\end{solution}

\end{parts}

\newpage

\question {\bf Smart contract programming:} {\em Matching Pennies}. This assignment will focus on writing your own smart contract to implement the \href{https://en.wikipedia.org/wiki/Matching_pennies}{Matching Pennies} game. The contract should allow two players (A, B) to play a game of Matching Pennies at any point in time. Each player picks a value of two options---for example, the options might be \{0, 1\}, \{`a’, ‘b’\}, \{True, False\}, etc. If both players pick the same value, the first player wins; if players pick different values, the second player wins. The winner gets 0.1 ETH as reward. After a game ends, two different players should be able to use your contract to play a new game.

\noindent {\bf Example:} Let A, B be two players who play the game, each with 1 ETH.  A picks 0 and B picks 0, so A wins. After the game ends, A’s balance is 1.1 ETH (perhaps minus some gas fees, if necessary).

You should implement the smart contract and deploy it on the course's Goerli Testnet.  Your contract should be as {\em secure}, {\em gas efficient}, and {\em fair} as possible. After deploying your contract, you should engage with other student's contract and play a game on his/her contract. Before you engage with a fellow student’s smart contract, you should evaluate their code and analyze its features in terms of security and fairness (refer to Lectures 07 and 08). You should provide:

\begin{parts}

\part[10] The code of your contract, together with a detailed description of the high-level decisions you made for the design of your contract,
including:  
\begin{itemize}
\item Who pays for the reward of the winner?
\item How is the reward sent to the winner?
\item How is it guaranteed that a player cannot cheat?
\item What data type/structure did you use for the pick options and why?
\end{itemize}

\begin{solution} %Uncomment and type your solution here
  Dejo acá el código fuente del contrato. 
  Tuve que formatearlo para que entre en el espacio horizontal del pdf.
  Si lo quieren leer un poco mejor, dejo el \href{https://github.com/pedrofuentes79/Fundamentos-y-Aplicaciones-de-Blockchain/blob/main/src/HW2_MatchingPennies.sol}{link del archivo en github} que uso para que lo puedan leer con linter.

  \begin{verbatim}
// SPDX-License-Identifier: GPL-3.0

// This makes sure we use a solidity version 
// that automatically handles overflow/underflow for me :D. 
// It reverts the operation if it overflows/underflows.
pragma solidity ^0.8.0;

contract MatchingPennies {
    address public playerA;
    address public playerB;
    bool public valueA;
    bool public revealedA;
    bool public valueB;
    bool public revealedB;
    bytes32 public commitmentA;
    bytes32 public commitmentB;
    uint256 public constant REWARD = 0.1 ether;
    uint256 public REVEAL_PERIOD = 1 hours;
    mapping(address => uint256) public balances;
    uint256 public constant MINIMUM_OPPONENT_WAIT = 5 minutes;
    uint256 public playerAWaitTime;
    uint256 public lastCommitmentTime;

    // Example:
    //   bytes32 commitment = keccak256(
    //     abi.encode(true, "mySecret123")
    //   );
    //   commitPlay(commitment) {value: 0.05 ether}
    function commitPlay(bytes32 commitment)
        public
        payable
    {
        require(
            msg.value == REWARD / 2,
            "Must pay half the reward to play"
        );

        // if A hasn't played yet
        if (playerA == address(0)) {
            playerA = msg.sender;
            commitmentA = commitment;
            // Tracks how long A is waiting for a player
            playerAWaitTime = block.timestamp;
        } else if (playerB == address(0)) {
            require(
                msg.sender != playerA,
                "Player A cannot play again"
            );
            playerB = msg.sender;
            commitmentB = commitment;
        } else {
            revert("Game already started");
        }

        // If both have committed, 
        // start the countdown for the reveal period
        if (playerA != address(0) && playerB != address(0))
        {
            lastCommitmentTime = block.timestamp;
        }
    }

    // Example:
    //   revealPlay(my_choice, "mySecret123"),
    //   where "mySecret123" is the secret used in commitPlay
    function revealPlay(
        bool value,
        string memory secret
    )
        public
    {
        require(
            playerA != address(0) && playerB != address(0),
            "Both players must have played"
        );

        bytes32 hash = keccak256(abi.encode(value, secret));

        if (msg.sender == playerA) {
            require(
                !revealedA, "Player A has already revealed"
            );
            require(
                hash == commitmentA,
                "Commitment does not match"
            );
            valueA = value;
            revealedA = true;
        } else if (msg.sender == playerB) {
            require(
                !revealedB, "Player B has already revealed"
            );
            require(
                hash == commitmentB,
                "Commitment does not match"
            );
            valueB = value;
            revealedB = true;
        } else {
            revert("Not a player");
        }

        if (revealedA && revealedB) {
            determineWinner();
        }
    }

    function determineWinner() private {
        require(
            revealedA && revealedB,
            "Both players must have revealed"
        );

        address winner;
        if (valueA == valueB) {
            winner = playerA;
        } else {
            winner = playerB;
        }
        // have the winner be able to retrieve the funds later.
        balances[winner] += REWARD;
        cleanUp();
    }

    function getBalance() public view returns (uint256) {
        return balances[msg.sender];
    }

    function withdraw() public {
        uint256 balance = balances[msg.sender];
        require(balance > 0, "No balance to withdraw");
        balances[msg.sender] = 0;
        payable(msg.sender).transfer(balance);
    }

    function cleanUp() private {
        playerA = address(0);
        playerB = address(0);
        revealedA = false;
        revealedB = false;
        valueA = false;
        valueB = false;
        commitmentA = bytes32(0);
        commitmentB = bytes32(0);
        lastCommitmentTime = 0;
    }

    function claimWinByTimeOut() public {
        require(
            playerA != address(0) && playerB != address(0),
            "Both players must have played"
        );
        require(
            block.timestamp - lastCommitmentTime
                > REVEAL_PERIOD,
            "Reveal period not over"
        );
        require(
            !(revealedA && revealedB),
            "Game already finished"
        );

        // Case 1: No one revealed
        if (!revealedA && !revealedB) {
            // we give the funds back to both players
            balances[playerA] += REWARD / 2;
            balances[playerB] += REWARD / 2;
        }
        // Case 2: Only one player revealed.
        // we give all the funds to the only one who revealed
        else {
            address winner;
            if (revealedA) {
                winner = playerA;
            } else {
                winner = playerB;
            }
            balances[winner] += REWARD;
        }
        cleanUp();
    }

    // This allows player A to get back their funds 
    // if no one played them yet
    // This can be done instantly, 
    // not necessary to wait for the reveal period.
    function forfeitIfNoOnePlayed() public {
        require(
            msg.sender == playerA,
            "Only player A can forfeit"
        );
        require(
            playerB == address(0),
            "Player B has already played"
        );
        require(
            block.timestamp - playerAWaitTime
                > MINIMUM_OPPONENT_WAIT,
            "Forfeit not allowed yet"
        );
        balances[playerA] += REWARD / 2;
        cleanUp();
    }
}


  \end{verbatim}
  
  Explicación del contrato:

  El premio del ganador lo paga quien perdió. Es decir, ambos jugadores (A y B) pagan la mitad del premio al participar (en el método $\texttt{commitPlay}$ )
  Para este método los participantes deben enviar un commitment de su elección. Este commitment lo generan con 
  
  $\texttt{bytes32 commitment = keccak256(abi.encode(true, "mySecret123"));}$.

  Junto con el commitment, deben enviar exactamente 0.05 ether.

  Luego, una vez que el segundo jugador ya jugó, el primer jugador puede revelar su elección con $\texttt{revealPlay(my\_choice, "mySecret123")}$.

  Ahí se tienen que enviar los mismos valores de $\texttt{my\_choice}$ y $\texttt{"mySecret123"}$ que se usaron antes.
  Una vez que $B$ revela su elección, se determina el ganador, que puede retirar sus fondos usando el método $\texttt{withdraw}$.

  Se garantiza que un jugador no pueda hacer trampas por los siguientes motivos:
  \begin{itemize}
    \item Ningún jugador puede ver la elección del otro jugador, ya que lo que se envía es un commitment. Si no fuese así, el segundo jugador siempre estaría en ventaja, pues podría ver en la blockchain la elección de $A$ y elegir lo opuesto para ganar.
    \item Las elecciones no pueden reveladas si ambos jugadores no commitearon aún.
    \item Se garantiza que un jugador no pueda cambiar su elección una vez que envió el commit, pues no podría haber obtenido el mismo hash con otra elección.
    \item Si el jugador $A$ no encontró a nadie en el tiempo especificado, puede retirar sus fondos usando el método $\texttt{forfeitIfNoOnePlayed}$.
    \item Si el jugador $A$ revela su elección y $B$ no lo hizo (o viceversa), y pasó mas de una hora desde que los dos jugadores hicieron commit, el que haya revelado gana.
      \subitem Esto es para evitar que un jugador revele su elección (que se haga visible en la blockchain por medio de la transacción de la función $\texttt{revealPlay}$) y el otro jugador no revele la suya (dado que sabe que va a perder, pues leyó lo de $A$). En este caso, quien ya reveló puede reclamar el premio (si pasó una hora)
  \end{itemize}

  Para la elección usé un booleano, ya que me pareció lo mas simple dado que solo hay dos opciones. 
  Guardé otro booleano para cada una de las elecciones, que me dice si el valor que tengo guardado en $\texttt{valueA}$ o $\texttt{valueB}$ fue efectivamente revelado.
  
  Otra opción podría haber siendo guardar un entero, con valores $0$, $1$ y $2$. El 0 representaría que no se ha revelado la opción aún.
  El 1 sería true y el 2 sería false.
\end{solution}

\part[5] A detailed gas consumption evaluation of your implementation, including: 
\begin{itemize}
\item The cost of deploying and interacting with your contract.
\item Whether your contract is {\em fair} to both players, including whether one player has to
pay more gas than the other and why.
\item Techniques to make your contract more cost efficient and/or fair.
\end{itemize}

\begin{solution} %Uncomment and type your solution here
  El costo de \href{https://sepolia.etherscan.io/tx/0x17fe5b5a1efd20dbfdb45436f2e9ebc09771e578830e30f4d8ad282e409a9e1a}{deploy} fue de 0.00279 ETH (desde Remix).

  Los costos de las transacciones fueron:
\begin{itemize}
  \item $A$ hace commit: 92,133 gas 
  \item $B$ hace commit: 92,473 gas 
  \item $B$ hace reveal: 36,338 gas 
  \item $A$ hace reveal: 75,156 gas 
  \item $A$ hace withdraw: 35,738 gas 
\end{itemize}
  El contrato $\textbf{no es justo}$ para ambos jugadores. Esto es porque, entre otros desbalances, el último en revelar tiene que pagar por la ejecución de la función
  $\texttt{deterimneWinner}$, y luego ejecuta $\texttt{cleanUp}$ para retornar al estado inicial. Esto se refleja en los costos de las transacciones dados.

  Para hacerlo mas eficiente, se podrían usar menos variables, los timestamps se podrían guardar en variables de tamaño más chico.
  Quizás para hacerlo más justo se podría hacer el cleanUp una vez que alguien quiera empezar un nuevo juego.

\end{solution}

\part[5] A thorough list of potential hazards and vulnerabilities that {\em may} occur in your contract; provide a detailed analysis of the security mechanisms you use to mitigate such hazards.

\begin{solution} %Uncomment and type your solution here
  \begin{itemize}
    \item Fondos bloqueados: Si $A$ y $B$ jugaron, pero $B$ no reveló su elección y $A$ sí lo hizo, $A$ puede reclamar su recompensa si pasó más de una hora desde que $B$ hizo commit (y no reveló su elección).
    \item Reentrancy: En el único momento en el que se sale de este contrato es en la función $\texttt{withdraw}$, cuando un jugador quiere retirar sus fondos. Y la llamada a $\texttt{transfer}$ se hace luego de haber actualizado el estado. 
    \item Frontrunning: Los jugadores tienen que hacer un commit con el hash de su elección + un secreto que elijan, y luego hacen el reveal con los mismos valores. 
         De este modo, nos aseguramos de que si alguien ve la transacción inicial (el commit) en la blockchain, no puede saber el valor de antemano. Si fuese así, el juego sería trivial para el jugador B
  \end{itemize}

\end{solution}

\part[5] A description of your analysis of your fellow student’s contract (along with relative code snippets of their contract, where needed for readability), including: 
\begin{itemize}
\item Any vulnerabilities discovered?
\item How could a player exploit these vulnerabilities to win the game?
\end{itemize}

\begin{solution} %Uncomment and type your solution here
    El código que estoy analizando es el de Juan D'Elia.

    \begin{verbatim}
// SPDX-License-Identifier: UNLICENSED
pragma solidity ^0.8.28;

contract MatchingPennies {
    struct Player {
        // si hago playersMapping[add] 
        // con una address que no esta, present sera 0
        bool present; 
        bool guessed;
        bool revealed;
        bool choice;
        bytes32 secret;
        uint256 balance; //balance que tiene para hacer withdraw
    }

    bool player1Ready;
    bool player2Ready;

    address player1Address;
    address player2Address;

    // si pasa cierto tiempo cualquiera puede resetear el juego
    uint256 public deadline;

    mapping(address => Player) public playersMapping;

    function updateSecret(
        address player,
        bytes32 _secret
    )
        private
    {
        require(
            !playersMapping[player].guessed,
            "No podes adivinar dos veces!"
        );
        playersMapping[player].secret = _secret;
        playersMapping[player].guessed = true;
    }

    function updateBalance(address player) private {
        playersMapping[player].balance += 0.2 ether;
    }

    function guess(bytes32 secret) public payable {
        require(
            playersMapping[msg.sender].present,
            "Solo pueden adivinar los jugadores de este juego"
        );
        require(
            msg.value == 0.1 ether,
            "Para jugar se debe enviar 0.1 ether"
        );
        require(
            player1Ready && player2Ready,
            "Ambos deben estar listos para empezar a jugar"
        );

        updateSecret(msg.sender, secret);
    }

    function resetPlayer(address _playerAddress) private {
        playersMapping[_playerAddress].present = false;
        playersMapping[_playerAddress].guessed = false;
        playersMapping[_playerAddress].revealed = false;
        playersMapping[_playerAddress].choice = false;
        playersMapping[_playerAddress].secret = 0;
    }

    function revealSecret(bool _choice, bytes32 secret) public {
        require(
            playersMapping[msg.sender].present,
            "Solo puede revelar su eleccion un jugador del juego"
        );
        require(
            playersMapping[msg.sender].secret
                == keccak256(abi.encodePacked(_choice, secret)),
            "El secreto y choice deben ser los originales"
        );

        playersMapping[msg.sender].revealed = true;
        playersMapping[msg.sender].choice = _choice;

        //si ambos ya revelaron 
        // puedo actualizar el balance del ganador
        if (
            playersMapping[player1Address].revealed
                && playersMapping[player2Address].revealed
        ) {
            bool choice1 = playersMapping[player1Address].choice;
            bool choice2 = playersMapping[player2Address].choice;
            if (choice1 == choice2) {
                updateBalance(player1Address);
            } else {
                updateBalance(player2Address);
            }
            //resetear a los jugadores 
            // y "desconectarlos" del juego
            resetPlayer(player1Address);
            resetPlayer(player2Address);
            player1Ready = false;
            player2Ready = false;
        }
    }

    function withdraw() public {
        require(
            playersMapping[msg.sender].balance > 0,
            "El balance debe ser mayor a cero"
        );

        uint256 balance = playersMapping[msg.sender].balance;
        playersMapping[msg.sender].balance = 0;
        payable(msg.sender).transfer(balance);
    }

    function joinGame() public {
        require(
            !player1Ready || !player2Ready,
            "Ya hay dos jugadores"
        );

        if (!player1Ready) {
            player1Address = msg.sender;
            playersMapping[player1Address].present = true;
            player1Ready = true;
        } else if (!player2Ready) {
            player2Address = msg.sender;
            playersMapping[player2Address].present = true;
            player2Ready = true;
            //cuando se une el segundo jugador pongo un deadline
            deadline = block.timestamp + 24 hours;
        }
    }
    //para evitar
    // 1. Que un jugador no deje que juegue nadie
    // 2. que cuando un jugador revele, 
        // el otro no lo haga porque sabe que perdio

    function restartGame() public {
        require(
            block.timestamp > deadline, "Deben pasar 24 horas"
        );
        require(
            player1Ready && player2Ready,
            "No hay ningun juego en curso"
        );
        //si algun jugador ya revelo y el otro hizo guess  
        // pero no revelo damos como ganador al que revelo
        if (playersMapping[player1Address].revealed) {
            // gana por abandono  1
            updateBalance(player1Address);
        } else if (playersMapping[player2Address].revealed) {
            // gana por abandono  2
            updateBalance(player2Address);
        }
        //si uno guessea y el otro no al pasarse el limite
        //  tiene la oportunidad de recuperar sus fondos
        else if (
            playersMapping[player1Address].guessed
                && !playersMapping[player2Address].guessed
        ) {
            playersMapping[player1Address].balance += 0.1 ether;
        } else if (
            playersMapping[player2Address].guessed
                && !playersMapping[player1Address].guessed
        ) {
            playersMapping[player2Address].balance += 0.1 ether;
        }

        //si ambos guessean pero ninguno revela o ninguno guessea 
        // el contrato se queda los fondos para penalizarlos

        resetPlayer(player1Address);
        resetPlayer(player2Address);
        player1Ready = false;
        player2Ready = false;
    }
}

    \end{verbatim}

    La lógica es bastante similar a la mía. 
    
    Un potencial detalle de seguridad podría ser de DoS. En el caso de que alguien solamente se una (sin pagar ningún costo, ya que el pago se hace en la funcion $\texttt{guess}$) y luego no juegue, podría hacer que,
    si otros jugadores quieren jugar, no puedan hacerlo, y deberá esperar hasta el deadline para poder reiniciar el juego y obtener de vuelta sus fondos.
    Sería una vulnerabilidad DoS ya que el contrato se vuelve inusable para otros jugadores hasta que pase el deadline.

    De todos modos, no es posible ganar el juego solamente con esto, lo único que hace es que un atacante pueda molestar a los que realmente quieren usar el contrato para jugar.
    
\end{solution}

\part[5] The transaction history of an execution of a game on your contract.

\begin{solution} %Uncomment and type your solution here


    - Jugador $A$ hace commit: \href{https://sepolia.etherscan.io/tx/0x1389d4e7496b2a47eefe4d1a1f3f9ae19323240f17465edf7ce603a77ccbebaf}{link}

    - Jugador $B$ hace commit: \href{https://sepolia.etherscan.io/tx/0x660cddd2aabb12fe5d679b687caa65227b99b4b5eda9f4fedc742af40a4652ea}{link}

    - Jugador $B$ hace reveal: \href{https://sepolia.etherscan.io/tx/0x06f5dc55924f04d033ccb4258754855af75550d6128c475f4cc7c993870a26ad}{link}

    - Jugador $A$ hace reveal: \href{https://sepolia.etherscan.io/tx/0xc5f6d423d6a042c85a5bec9638c7f4d9114e1020bd527b43973f4c14700c88fd}{link}

    - Jugador $A$ hace withdraw: \href{https://sepolia.etherscan.io/tx/0x09f1f95ec3edef7e1b57208f6060f4a747f45ab0ca4db598c7ea19eaa5de75f1}{link}


    $A$ ganó ya que había elegido $true$ y $B$ había elegido $true$. Luego, $A$ pudo retirar un total de 0.1 ether.
\end{solution}

\end{parts}

\newpage

\newpage

~\\

\end{questions}

\end{document}
